\clearpage
\VerdierPhysicalPageAnchor{0181}%
\begingroup\small\setlength{\parskip}{0pt}\setlength{\emergencystretch}{1.5em}
\setlength{\abovedisplayskip}{4pt plus 2pt minus 2pt}%
\setlength{\belowdisplayskip}{4pt plus 2pt minus 2pt}%

\VerdierSemanticAnchor{ch3d-1-2-10-proof-cont}
Le complexe $X(0,+\infty)$ (démonstration de (1.2.1)) est alors
quasi-isomorphe à $X$, et le morphisme canonique
$H^0(X(0,+\infty))\longrightarrow X(0,+\infty)$ est un
quasi-isomorphisme. Ce qui montre que $X$ est isomorphe dans
$D(\mathcal A)$ à $H^0(X)$.

Démontrons b). La deuxième égalité est mise pour mémoire et résulte de ce
que $\mathcal A$ est une sous-catégorie pleine de $D(\mathcal A)$. Un
morphisme dans $D(\mathcal A)$ de $X[-n]$ dans $Y$ est de la forme
$Q(m)Q(s)^{-1}$, où $m\colon Z^\bullet\longrightarrow Y$ est un morphisme
de $K(\mathcal A)$ et $s\colon Z^\bullet\longrightarrow X[-n]$ est un
quasi-isomorphisme de $K(\mathcal A)$. Le morphisme
$i\colon Z^\bullet(-\infty,n+1)\longrightarrow Z^\bullet$ (démonstration de
(1.2.1)) est alors un quasi-isomorphisme. Par suite, on a :
\[
Q(m)Q(s)^{-1}=Q(m)Q(i)Q(i)^{-1}Q(s)^{-1}=Q(mi)Q(si)^{-1}.
\]
Lorsque $n$ est strictement négatif, $mi=0$. Par suite,
$Q(m)Q(s)^{-1}=0$.

\VerdierSemanticAnchor{ch3d-1-2-11}
\medskip\noindent\textbf{Définition 1.2.11.}\quad
On désigne par $\mathcal A^{\operatorname{Ext}}$ la sous-catégorie pleine de
$D(\mathcal A)$ définie par les translatés des objets de $\mathcal A$.

La catégorie $\mathcal A^{\operatorname{Ext}}$ est donc (contrairement à la
sous-catégorie $\mathcal A$) une sous-$\mathbb Z$-catégorie de
$D(\mathcal A)$. Un objet de $D(\mathcal A)$ est isomorphe à un objet de
$\mathcal A^{\operatorname{Ext}}$ si et seulement s'il ne possède qu'un seul
objet de cohomologie non nul.

\VerdierSemanticAnchor{ch3d-1-3}
\subsection{Le \texorpdfstring{$\delta$}{delta}-foncteur canonique.}

\VerdierSemanticAnchor{ch3d-1-3-1}
\noindent\textbf{1.3.1.}\quad
Soit $\mathcal A$ une catégorie abélienne et $\mathcal D$ une catégorie
triangulée. Un $\delta$-foncteur de $\mathcal A$ dans $\mathcal D$ est un
objet constitué par :

\noindent 1) Un foncteur $G\colon\mathcal A\longrightarrow\mathcal D$.

\noindent 2) La donnée, pour toute suite exacte
$S=0\longrightarrow X\longrightarrow Y\longrightarrow Z\longrightarrow0$
de $\mathcal A$, d'un morphisme de $\mathcal D$ :
\[
\delta(S)\colon G(Z)\longrightarrow G(X)[1].
\]
Ces données étant soumises aux axiomes suivants :

\noindent a) Pour tout morphisme de suites exactes :

\VerdierSemanticAnchor{ch3d-1-3-1-exact-sequences}
\begin{center}
\resizebox{0.88\textwidth}{!}{%
$\begin{tikzcd}[ampersand replacement=\&,column sep=2.6em,row sep=2.2em]
S={}\&0\arrow[r]\&X\arrow[r]\arrow[d,"u"']\&Y\arrow[r]\arrow[d,"v"']\&Z\arrow[r]\arrow[d,"w"']\&0\\
S'={}\&0\arrow[r]\&X'\arrow[r]\&Y'\arrow[r]\&Z'\arrow[r]\&0
\end{tikzcd}$%
}
\end{center}
\endgroup
